(1) 记主质量位移 \(u_n\)，内部位移 \(v_n\)。方程
\[
M\ddot u_n=K(u_{n+1}+u_{n-1}-2u_n)+k(v_n-u_n),\quad m\ddot v_n=k(u_n-v_n).
\]
Bloch代入给 \(V/U=k/(k-m\omega^2)\)，主方程化为
\[
4K\sin^2(qa/2)=\omega^2 M_{\rm eff}(\omega),\quad M_{\rm eff}=M+\frac{mk}{k-m\omega^2}.
\]
(2) 在 \(\omega_0<\omega<\omega_+\) 时，分母为负且 \(M_{\rm eff}<0\)（其零点在 \(\omega_+\)）。左边对实 \(q\) 非负，右边严格负，故无传播Bloch模。此禁带由胞内共振惯量变号产生，不依赖晶格半波长Bragg反射。
(3) 低频 \(M_{\rm eff}\to M+m\)，\(\sin(qa/2)\sim qa/2\)，故 \(\omega\simeq a\sqrt{K/(M+m)}|q|\)。在 \(q=0\)，色散行列式的两根为 \(\omega=0\) 和 \(\omega_+\)。\(V/U=k/(k-m\omega^2)\) 在 \(\omega\to\omega_0\) 形式发散；对非零驱动该点内部共振强烈，必须考虑阻尼或非线性，而非把无穷大位移视作物理结果。
(4) 每胞瞬时能量
\[
E_n=\frac M2\dot u_n^2+\frac m2\dot v_n^2+\frac K2(u_{n+1}-u_n)^2+\frac k2(v_n-u_n)^2\ge0.
\]
仅整体平移的零频模能使所有弹簧不变形。动态消元的 \(M_{\rm eff}<0\) 是频域响应符号，不是任何单颗粒惯性质量或真实动能为负。加入小阻尼使有效质量复数化，禁带变成强衰减带；大振幅使线性弹簧近似失效。
